\chapter{Identificación del problema}

\section{Flujo productivo del HTP}

El HTP del A320 es un conjunto aeronáutico de alta criticidad. Su fabricación y montaje requieren útiles específicos, documentación, control dimensional, trazabilidad de materiales, inspecciones y coordinación entre estaciones. En el entorno de sección 19.1, la operación se relaciona con el área posterior del fuselaje y las interfaces de cola, por lo que la disponibilidad correcta de kits, herramientas y consumibles condiciona la continuidad del trabajo.

De forma simplificada, el flujo operativo incluye preparación de kits, traslado a estación, identificación de material, uso de útiles y herramientas, inspección, retirada de elementos usados, registro y cierre documental. Cada paso puede funcionar correctamente de manera aislada y, aun así, generar esperas si la coordinación logística no está sincronizada con el ritmo de producción.

\section{Problemas típicos detectados}

\begin{itemize}
  \item \textbf{Desplazamientos de operarios.} Parte del tiempo se dedica a buscar, recoger o devolver útiles y consumibles en lugar de ejecutar tareas de valor directo.
  \item \textbf{Esperas por kits o herramientas.} Un kit incompleto o una herramienta no localizada bloquea la secuencia aunque la estación este disponible.
  \item \textbf{Riesgo FOD.} Restos, tornilleria, embalajes o elementos no controlados pueden generar incidencias de calidad y seguridad.
  \item \textbf{Trazabilidad fina insuficiente.} Saber que un lote existe no siempre equivale a conocer en tiempo real donde esta, quien lo recibio y para que orden se uso.
  \item \textbf{Coordinación entre estaciones.} Las prioridades cambian durante el turno y el suministro manual puede reaccionar tarde.
  \item \textbf{Escalado de capacidad.} Pasar de 40 HTP al mes a 80 HTP al mes o más exige reducir fricciones operativas; no basta con duplicar recursos sin mejorar el flujo.
\end{itemize}

\section{Oportunidad de automatización}

La oportunidad está en automatizar el movimiento y el registro, no en reemplazar decisiones de fabricación. Un AMR puede ejecutar misiones recurrentes, entregar kits a una ventana horaria, verificar identificaciones RFID/QR, registrar confirmaciones, inspeccionar zonas definidas y devolver útiles usados. El valor proviene de estabilizar el flujo y reducir incertidumbre, especialmente cuando la planta opera cerca de su límite de capacidad.
